\documentclass[11pt, a4paper]{article}

% --- Packages ---
\usepackage[utf8]{inputenc}
\usepackage[T1]{fontenc}
\usepackage{geometry}
\geometry{top=2.5cm, bottom=2.5cm, left=2.5cm, right=2.5cm}
\usepackage{mathpazo} % Palatino font (exceptional readability)
\usepackage[dvipsnames]{xcolor}
\usepackage{hyperref}
\usepackage{listings} % For code snippets
\usepackage{tcolorbox} % For highlighted boxes
\usepackage{graphicx}
\usepackage{enumitem}
\usepackage{fancyhdr}
\usepackage{booktabs} % For professional tables

% --- Styling ---
\definecolor{codegray}{rgb}{0.95,0.95,0.95}
\definecolor{codeblue}{rgb}{0.1,0.1,0.5}
\definecolor{codegreen}{rgb}{0,0.5,0}

\lstset{
    backgroundcolor=\color{codegray},
    commentstyle=\color{codegreen},
    keywordstyle=\color{magenta},
    numberstyle=\tiny\color{gray},
    stringstyle=\color{codeblue},
    basicstyle=\ttfamily\small,
    breakatwhitespace=false,
    breaklines=true,
    captionpos=b,
    keepspaces=true,
    numbers=left,
    numbersep=5pt,
    showspaces=false,
    showstringspaces=false,
    showtabs=false,
    tabsize=2,
    frame=single
}

\pagestyle{fancy}
\fancyhf{}
\rhead{\textbf{ClusterLLM-Local}}
\lhead{Day 1: Infrastructure \& Philosophy}
\rfoot{\thepage}

% --- Document Info ---
\title{\textbf{\Huge Day 1 Technical Report}\\ \Large Reproducibility, Dependency Management (Poetry), and Data Ops (DVC)}
\author{\textbf{Hamady Gackou} \\ Master Data Science \& AI}
\date{\today}

\begin{document}

\maketitle
\thispagestyle{empty}

\begin{abstract}
    This document details the foundational engineering choices made on Day 1 of the \textit{ClusterLLM-Local} project. Beyond setting up the environment on a GCP L4 GPU instance, this report analyzes the strategic decisions to prioritize \textbf{Reproducibility} over speed. It provides a technical justification for choosing \textbf{Poetry} over standard \texttt{pip/venv} for dependency management and \textbf{DVC} for data versioning. The day concluded with the successful deployment of a "Tracer Bullet" pipeline using the \textit{Bank77} dataset.
\end{abstract}

\tableofcontents
\vspace{1cm}
\hrule
\vspace{1cm}

\section{The Core Philosophy: Reproducibility}

Before writing code, we defined the primary constraint of this research: \textbf{Reproducibility}.

\subsection{What is Reproducibility?}
In the context of Machine Learning, reproducibility is the ability for an independent researcher (or a Jury member) to recreate the exact same results (metrics, graphs, models) given the same input data and code. 

\subsection{Why is it critical for this project?}
\begin{itemize}
    \item \textbf{Scientific Rigor (Master's Defense):} If the code works "only on my machine" because of a hidden manual configuration, the scientific contribution is null.
    \item \textbf{The "Dependency Hell" Risk:} In 6 months, libraries like \texttt{transformers} or \texttt{torch} will update. Without strict locking, the code might break or produce different results.
    \item \textbf{PhD Track Requirement:} Top-tier conferences (ICLR, NeurIPS) require a "Reproducibility Checklist". Setting this up on Day 1 ensures compliance.
\end{itemize}

\section{Dependency Management: Why Poetry?}

To ensure reproducibility, we rejected the standard \texttt{venv + pip} workflow in favor of \textbf{Poetry}.

\subsection{The Limitation of Standard Virtual Environments}
A standard Python workflow uses `virtualenv` and a `requirements.txt` file generated by `pip freeze`.
\begin{itemize}
    \item \textbf{The Problem:} `pip` installs packages sequentially. If Package A needs `numpy<2.0` and Package B needs `numpy>=2.0`, pip might install conflicting versions without warning, leading to runtime crashes.
    \item \textbf{No Locking:} A `requirements.txt` often lists top-level packages (e.g., `pandas`). It rarely locks the hundreds of sub-dependencies. A user installing it tomorrow might get a different sub-dependency version than the one I used today.
\end{itemize}

\subsection{The Poetry Solution}
Poetry is a dependency manager and packaging tool that solves these issues. 

\begin{enumerate}
    \item \textbf{Deterministic Resolution:} Before installing anything, Poetry solves the mathematical graph of dependencies to ensure no conflicts exist.
    \item \textbf{The Lock File (\texttt{poetry.lock}):} This is the "Source of Truth". It records the exact SHA-256 hash of every single file installed. 
    \item \textbf{Separation of Concerns:} It separates \textit{Production} dependencies (PyTorch) from \textit{Development} dependencies (PyTest, Black) in the `pyproject.toml` file.
\end{enumerate}

\begin{tcolorbox}[colback=green!5!white,colframe=green!75!black,title=Analogy]
    \textbf{Standard Venv:} Throwing ingredients into a pot and hoping the soup tastes good.\\
    \textbf{Poetry:} A strict recipe with exact measurements and brand names for every ingredient, checked by a chemist before cooking.
\end{tcolorbox}

\section{Data Operations: Data Version Control (DVC)}

Just as Poetry locks the code environment, DVC locks the data.

\subsection{The Problem with Git and Data}
Git computes differences between lines of text. It cannot handle binary large objects (BLOBs) like CSVs or Model Weights. Storing a 500MB dataset in Git slows down the repository and hits storage limits.

\subsection{How DVC Works}
DVC acts as a pointer system (like a hyperlink to a file). 

\begin{enumerate}
    \item \textbf{The Workspace:} We put the actual data in `data/raw/bank77.csv`.
    \item \textbf{The Hashing:} \texttt{dvc add} calculates the MD5 hash of the file.
    \item \textbf{The Pointer:} DVC replaces the data with a small text file: `bank77.csv.dvc`.
    \item \textbf{The Remote:} The actual heavy file is pushed to Google Drive (our "Object Storage").
\end{enumerate}

\begin{table}[h]
\centering
\begin{tabular}{@{}lll@{}}
\toprule
\textbf{Feature} & \textbf{Git} & \textbf{DVC} \\ \midrule
Handles & Code (Text) & Data \& Models (Binary) \\
Storage & GitHub Servers & Google Drive / S3 / GCP \\
Tracking Method & Line-by-line Diff & MD5 Hash \\
Repo Size & Small (<100MB) & Small (Contains only pointers) \\ \bottomrule
\end{tabular}
\caption{Comparison of Git vs. DVC roles in the project.}
\end{table}

\section{Implementation Log (Day 1)}

\subsection{1. Environment Setup}
We initialized the deterministic environment.
\begin{lstlisting}[language=bash, caption=Poetry Setup]
# Core dependencies for the research paper (CUDA 11.8 compatible)
poetry add torch torchvision torchaudio --index-url https://download.pytorch.org/whl/cu118
# MLOps Stack
poetry add dvc dvc-gdrive mlflow hydra-core
\end{lstlisting}

\subsection{2. Connecting Storage}
We linked the local GPU instance to Google Drive.
\begin{lstlisting}[language=bash, caption=DVC Remote Setup]
dvc init
dvc remote add -d myremote gdrive://<GDRIVE_ID>
git commit -m "chore: setup dvc remote"
\end{lstlisting}

\subsection{3. The Tracer Bullet (Bank77)}
To validate the infrastructure without incurring the cost of processing all 14 datasets, we implemented a single dataset pipeline.

\begin{lstlisting}[language=bash, caption=Workflow Execution]
# 1. Download Data
poetry run python src/data/loader.py

# 2. Version Control Data (DVC)
dvc add data/raw/bank77_train.csv
dvc push  # Sends to Google Drive

# 3. Version Control Code (Git)
git add data/raw/*.dvc
git commit -m "data: add bank77 dataset"
git push
\end{lstlisting}

\section{Conclusion \& Roadmap}

We have successfully established a \textbf{Reproducible Infrastructure}.
\begin{itemize}
    \item Code is locked via \texttt{poetry.lock}.
    \item Data is locked via \texttt{.dvc} files.
    \item History is tracked via \texttt{git}.
\end{itemize}

\textbf{Next Step (Day 2):} Containerization. We will wrap this entire ecosystem into a \textbf{Docker} container to guarantee that the operating system libraries (CUDA drivers) are also reproducible.

\end{document}